% Section 2 -- The DLC Calculus
% Target: 3 pages.
% Goal: present the syntax, judgment forms, and key typing rules.

\subsection{Syntactic categories}

DLC's grammar is parameterized over principal identifiers ($h$ ranging
over SPIFFE-ID hashes), atom names ($a$), action identifiers ($\alpha$),
time bounds ($\tau$), algorithm tags ($\gamma$), and signature byte
strings ($\sigma$).

\begin{figure}[h]
\small
\begin{align*}
p, q, r &::= h \mid p \wedge q \mid p \vee q \mid p \sqcap q
& \text{(principals)} \\
\ell &::= \langle i_1, \ldots, i_n \rangle
& \text{(IFC labels)} \\
O &::= \top \mid \bot \mid \mathsf{act\_of}(p, \alpha) \mid \mathsf{within}(\tau)
\mid O_1 \otimes O_2 \mid O_1 \multimap O_2
& \text{(obligations)} \\
\varphi, \psi &::= a \mid \top \mid \bot \mid \varphi \supset \psi
\mid \varphi \wedge \psi \mid \varphi \vee \psi
\mid p~\mathsf{says}~\varphi \mid p \Rightarrow q
& \text{(propositions)} \\
&\quad \mid \varphi @ \ell \mid \Box_O \varphi \mid \Diamond_\tau \varphi
\mid \varphi \otimes \psi \mid \varphi \multimap \psi \\
M, N &::= x \mid \lambda x{:}\varphi.\, M \mid M\,N
\mid \langle M, \sigma \rangle_p \mid \mathsf{verify}_p(M, \sigma)
& \text{(proof terms)} \\
&\quad \mid \mathsf{delegate}(M, N) \mid \mathsf{attenuate}(M, \psi)
\mid \mathsf{discharge}(M, N) \\
&\quad \mid \mathsf{lift}_\ell(M) \mid \mathsf{declassify}_\ell(M, \pi)
\mid \mathsf{now}(\tau) \mid \mathsf{within}(M, \tau) \\
&\quad \mid \langle M, N \rangle \mid \pi_1 M \mid \pi_2 M \mid \mathsf{inl}_\psi M
\mid \mathsf{inr}_\varphi M \mid \mathsf{case}~M~\mathsf{of}~\mathsf{inl}~x \Rightarrow N \mid \mathsf{inr}~y \Rightarrow P \\
&\quad \mid M \otimes N \mid \mathsf{let}~x \otimes y = M~\mathsf{in}~N
\mid \mathsf{let}~\langle x \rangle_p = M~\mathsf{in}~N \mid \mathsf{sfExtract}(M)
\end{align*}
\caption{DLC syntactic categories.}
\label{fig:syntax}
\end{figure}

The operator $\sqcap$ is intentionally associative but not commutative:
$h \sqcap h' \sqcap h''$ reads ``$h''$ acting as $h'$ acting as $h$'', the
left-associative chain structure that RFC~8693 fudges as an opaque
transitive closure.

\subsection{Judgments}

DLC has four mutually-recursive judgments:

\begin{itemize}
\item $\Gamma \vdash M : \varphi$ -- standard logical typing.
\item $\Gamma \vdash M \rhd M'$ -- small-step reduction.
\item $\Gamma \vdash_K M : \varphi$ -- cryptographic typing under
  keyring $K$ (the rule replaces every \textsc{says-I} signature
  with a cryptographic verification step against $K$).
\item Membership and well-formedness conditions on $\Gamma$ and $K$.
\end{itemize}

The bridge between $\vdash$ and $\vdash_K$ is T2 (Section~\ref{sec:metatheory}).

\subsection{Key typing rules}

We present the rules carrying the contribution claims.  The full
system (28 \textsf{Deriv} rules over 22 term constructors) is in the
companion artifact.  Caveat: two rules are mechanized in weaker form
than presented here --- \textsc{attenuate} is degenerate ($\psi =
\varphi$; the propositional-implication side condition is deferred)
and \textsc{discharge}'s obligation evidence is a placeholder
(\texttt{atom 0}) pending the obligation-carrying term constructor
(see Section~\ref{sec:metatheory}, T4).

\paragraph{Affirmation introduction (\textsc{says-I})}

\begin{mathpar}
\inferrule[says-I]
  {\Gamma \vdash M : \varphi \\ \sigma = \mathsf{Sign}_{sk(p)}(\mathsf{canonical}(M))}
  {\Gamma \vdash \langle M, \sigma \rangle_p : p~\mathsf{says}~\varphi}
\end{mathpar}

The signature $\sigma$ is intrinsic to the proof term, not an
external annotation.  T2 (cryptographic correspondence) is the
load-bearing claim that this intrinsic shape coincides with the
operational signing/verification under $K$.

\paragraph{Delegation (\textsc{delegate})}

\begin{mathpar}
\inferrule[delegate]
  {\Gamma \vdash M : p~\mathsf{says}~(q \Rightarrow p)
   \\
   \Gamma \vdash N : q~\mathsf{says}~\varphi}
  {\Gamma \vdash \mathsf{delegate}(M, N) : (p \sqcap q)~\mathsf{says}~\varphi}
\end{mathpar}

\textbf{The chain-splicing-defeating rule.}  The variable $q$ is bound
in both premises; an attacker who has a $\mathsf{says}~(q \Rightarrow p)$
token and a $\mathsf{says}~\varphi$ token by a different principal $q' \neq q$
cannot apply this rule.  The same condition propagates to the symbolic
model (Section~\ref{sec:symbolic}) and the wire format (Section~\ref{sec:wire}).

\paragraph{Attenuation (\textsc{attenuate})}

\begin{mathpar}
\inferrule[attenuate]
  {\Gamma \vdash M : p~\mathsf{says}~\varphi
   \\
   \varphi \supset \psi
   \\
   \psi \leq_L \varphi}
  {\Gamma \vdash \mathsf{attenuate}(M, \psi) : p~\mathsf{says}~\psi}
\end{mathpar}

The side condition $\psi \leq_L \varphi$ requires that attenuation
respects the IFC lattice on labels---preventing privilege-escalation
attacks that look like ordinary attenuation but raise an integrity label.

\paragraph{Discharge (\textsc{discharge})}

\begin{mathpar}
\inferrule[discharge]
  {\Gamma \vdash M : \Box_O \varphi
   \\
   \Gamma \vdash N : O \text{ (linear)}}
  {\Gamma \vdash \mathsf{discharge}(M, N) : \varphi}
\end{mathpar}

The only elimination form for $\Box_O \varphi$.  The linearity of $N$'s
typing enforces single-use semantics: an obligation discharged once
cannot be discharged again, preserving the multiset accounting that
T4 (obligation soundness) tracks.

\subsection{Reduction}

DLC's reduction relation includes the standard $\beta$-reduction plus
calculus-specific redexes:

\begin{align*}
(\lambda x{:}\varphi.\,M)\,N &\rhd M[N/x] \\
\mathsf{delegate}(\langle M, \_ \rangle_p, \langle N, \sigma' \rangle_q)
&\rhd \langle N, \sigma' \rangle_{p \sqcap q} \\
\pi_1 \langle a, b \rangle &\rhd a \qquad \pi_2 \langle a, b \rangle \rhd b \\
\mathsf{case}\,(\mathsf{inl}_\psi\,a)\,\ldots &\rhd \mathsf{left}[a/x] \\
\mathsf{let}~x \otimes y = a \otimes b~\mathsf{in}~M &\rhd M[a/x][b/y]
\end{align*}

Subject reduction (T4's antecedent) is stated in Section~\ref{sec:metatheory}.
